% cap_implementacion.tex — Capítulo: Desarrollo e implementación
% FINARA — TT 2026-A062 | IPN ESCOM
%
% Contiene la documentación técnica de la aplicación móvil Flutter.
% Fuente: documentacion_modulos.tex (módulos principales de la app).
%
% Pendiente integrar:
%   - Documentación del backend (Node.js + Express + MySQL)

\chapter{Desarrollo e implementación}

% ══════════════════════════════════════════════════════════════
\section{Aplicación móvil}
% ══════════════════════════════════════════════════════════════

\subsection{Arquitectura general de la aplicación}

\subsubsection{Stack tecnológico}

\begin{longtable}{>{\bfseries\small}p{3.5cm} p{10cm}}
  \caption{Stack tecnológico de la aplicación móvil FINARA}
  \label{tab:stack-movil} \\
  \toprule
  Componente & Descripción \\
  \midrule
  \endfirsthead
  \toprule
  Componente & Descripción \\
  \midrule
  \endhead
  Flutter SDK & Framework principal para la aplicación móvil multiplataforma.
    Versión \textasciicircum 3.6.0. Target: Android. \\[4pt]
  Dart & Lenguaje de programación tipado para el frontend móvil. \\[4pt]
  Node.js + Express & Backend REST API. Maneja autenticación JWT, lógica de
    negocio, invocación del pipeline OCR y persistencia. \\[4pt]
  MySQL & Base de datos relacional con tablas para usuarios, categorías,
    tickets, productos, presupuestos, deudas y alertas. \\[4pt]
  Python 3 (Tesseract) & Script \archivo{LecturaTickets\_v2.py} invocado
    como subproceso por Node.js. Ejecuta el pipeline OCR completo. \\[4pt]
  Google Cloud Vision & API de visión por computadora usada como respaldo
    cuando la confianza de Tesseract es insuficiente. \\[4pt]
  Claude Haiku & Modelo de lenguaje usado para estructurar el texto extraído
    por Vision en JSON de productos. \\[4pt]
  Firebase Crashlytics & Monitoreo de errores y excepciones en producción. \\[4pt]
  SharedPreferences & Almacenamiento local del token JWT y estado de sesión. \\
  \bottomrule
\end{longtable}

\subsubsection{Estructura del proyecto Flutter}

El proyecto sigue una arquitectura \textbf{feature-first}: cada módulo
funcional reside en su propio directorio dentro de \archivo{lib/features/},
con sus propias pantallas, servicios y widgets. El núcleo compartido se
ubica en \archivo{lib/core/}.

\begin{lstlisting}[style=codigobase, language={},
    caption={Estructura de directorios del proyecto Flutter}]
lib/
|-- core/
|   |-- constants/
|   |   |-- app_colors.dart       (paleta global de colores)
|   |   `-- api_constants.dart    (URLs de endpoints)
|   |-- models/
|   |   `-- finance_models.dart   (modelos: Budget, Debt, Category...)
|   `-- services/
|       `-- api_service.dart      (cliente HTTP centralizado)
`-- features/
    |-- scan/          (escaneo OCR de tickets)
    |-- categories/    (gestion de categorias)
    |-- budgets/       (presupuestos mensuales)
    |-- debts/         (deudas e intereses)
    |-- transactions/  (gastos e ingresos)
    |-- home/          (pantalla principal)
    |-- auth/          (login, registro)
    |-- notifications/ (alertas y configuracion)
    `-- profile/       (perfil de usuario)
\end{lstlisting}

\subsubsection{Comunicación con el backend}

Toda la comunicación entre la aplicación móvil y el backend se realiza
mediante \textbf{API REST} sobre HTTPS. El cliente HTTP centralizado
\campo{ApiService} (singleton) gestiona:

\begin{itemize}
  \item Inyección automática del token JWT en el header
        \campo{Authorization}.
  \item Decodificación de respuestas en formato \campo{\{success, data\}}.
  \item Timeouts configurables: 30 segundos para peticiones normales,
        60 segundos para el endpoint OCR.
  \item Manejo de errores HTTP (400, 401, 403, 404, 409, 500).
  \item Subida de imágenes mediante \campo{multipart/form-data}.
\end{itemize}

La URL base del backend es:
\endpoint{https://documents.digitaltax-analytics.com/finara/api}

\subsubsection{Manejo de estado}

La aplicación utiliza el patrón \textbf{StatefulWidget + setState} de
Flutter sin gestores de estado globales. Cada pantalla administra su
propio estado local. La capa de servicios (\campo{*Service}) actúa como
intermediaria entre la UI y el backend, abstrayendo la lógica HTTP.

\begin{notabox}
  No se utilizan paquetes como Provider, Riverpod o Bloc. El estado es
  local a cada widget. La comunicación entre pantallas se realiza mediante
  valores de retorno en \campo{Navigator.push} (\campo{result == true})
  que disparan recargas en la pantalla padre.
\end{notabox}

\subsubsection{Paleta de colores del sistema}

\begin{table}[H]
  \centering
  \caption{Paleta de colores de la aplicación móvil FINARA}
  \label{tab:paleta-movil}
  \begin{tabular}{lll}
    \toprule
    \textbf{Nombre} & \textbf{Hex} & \textbf{Uso} \\
    \midrule
    \campo{prussianBlue} & \#0f172a & Color primario, textos, AppBar \\
    \campo{seaGreen}     & \#13853d & Ingresos, éxito, indicadores positivos \\
    \campo{mahoganyRed}  & \#b51818 & Gastos, errores, alertas críticas \\
    \campo{platinum}     & \#f1f5f9 & Fondo de pantallas \\
    \campo{white}        & \#ffffff & Superficies de tarjetas \\
    \bottomrule
  \end{tabular}
\end{table}


% ──────────────────────────────────────────────────────────────
\subsection{Módulo de escaneo de tickets (OCR)}
% ──────────────────────────────────────────────────────────────

\subsubsection{Descripción general}

El módulo de escaneo es el componente más diferenciador de FINARA. Permite
al usuario fotografiar uno o varios tickets de compra y extraer
automáticamente los productos, precios, tienda y total mediante un pipeline
de OCR de múltiples etapas. El resultado alimenta directamente el registro
de gastos y el control de presupuestos.

\subsubsection*{Flujo completo de operación}

\begin{enumerate}
  \item \textbf{Captura} (\campo{ScanTicketScreen}): El usuario activa la
    cámara trasera y captura entre 1 y 5 fotografías del ticket.
  \item \textbf{Previsualización} (\campo{ScanPreviewScreen}): Carrusel de
    imágenes con opción de eliminar, añadir más fotos o confirmar.
  \item \textbf{Procesamiento} (\campo{ScanProcessingScreen}): Las imágenes
    se suben al backend vía \campo{multipart/form-data}. El backend invoca
    el script Python OCR y devuelve el JSON estructurado.
  \item \textbf{Resultado} (\campo{ScanResultScreen}): El usuario revisa,
    edita o elimina los productos detectados antes de guardar el gasto.
\end{enumerate}

\subsubsection{Pantallas del módulo}

\subsubsection*{ScanTicketScreen — Captura de imágenes}

Utiliza el plugin \campo{camera} para inicializar la cámara trasera con
resolución máxima.

\begin{lstlisting}[style=dart,
    caption={Inicialización de cámara en ScanTicketScreen}]
// Inicializacion de camara con permiso y resolucion
Future<void> _initializeCamera() async {
  final cameras = await availableCameras();
  _controller = CameraController(
    cameras.first,
    ResolutionPreset.max,
    imageFormatGroup: ImageFormatGroup.jpeg,
  );
  await _controller.initialize();
}
\end{lstlisting}

Características relevantes:

\begin{itemize}
  \item Máximo de 5 imágenes por ticket para controlar el tamaño del
        payload.
  \item Visualización en tiempo real del encuadre con guía superpuesta
        (CustomPainter).
  \item Botón de flash (\campo{FlashMode.torch}) para ambientes oscuros.
  \item Galería inferior deslizable con las fotos capturadas y botón de
        eliminación individual.
\end{itemize}

\subsubsection*{ScanPreviewScreen — Revisión previa}

Muestra las imágenes capturadas en un \campo{PageView} con navegación
manual. El usuario puede:

\begin{itemize}
  \item Eliminar la imagen actual con confirmación.
  \item Regresar a la cámara para añadir más fotos (respeta el límite
        de 5).
  \item Confirmar el lote de imágenes y pasar al procesamiento.
\end{itemize}

\subsubsection*{ScanProcessingScreen — Ejecución OCR}

Es la pantalla de espera durante el procesamiento. Muestra 6 pasos
animados que orientan al usuario sobre el progreso:

\begin{enumerate}
  \item Preparando imágenes
  \item Subiendo al servidor
  \item Ejecutando OCR
  \item Extrayendo texto
  \item Identificando productos
  \item Calculando totales
\end{enumerate}

La petición al backend se realiza mediante \campo{TicketsService.scanTicket()}:

\begin{lstlisting}[style=dart,
    caption={Llamada al endpoint de escaneo OCR}]
// Subida multipart al endpoint POST /api/tickets/scan
Future<Map<String, dynamic>> scanTicket(List<String> imagePaths) async {
  return await _apiService.postMultipart(
    ApiConstants.ticketsScan,
    imagePaths,
  );
}
\end{lstlisting}

Si el OCR falla, la pantalla presenta un estado de error con botones para
reintentar o cancelar (regresa al inicio sin persistir nada).

\subsubsection*{ScanResultScreen — Revisión y guardado}

Muestra los productos extraídos en una lista editable. El usuario puede:

\begin{itemize}
  \item Editar nombre, precio y categoría de cada producto.
  \item Eliminar productos no deseados.
  \item Añadir productos manualmente que no fueron detectados.
  \item Cambiar la fecha de la compra con un \campo{DatePicker}.
  \item Ver la confianza OCR del escaneo (indicador verde/naranja).
  \item Guardar el gasto completo.
\end{itemize}

\begin{lstlisting}[style=dart,
    caption={Guardado del gasto con productos extraídos por OCR}]
// Guardado del gasto con productos extraidos
Future<void> _saveExpense() async {
  await _expensesService.createWithProducts(
    idTicket:    _ticketId,
    nombreGasto: _expenseNameController.text,
    fecha:       _selectedDate,
    productos:   _extractedItems.map((item) => {
      'nombre_producto': item['name'],
      'precio':          item['price'],
      'id_categoria':    item['id_categoria'],
    }).toList(),
  );
}
\end{lstlisting}

\subsubsection{Pipeline OCR en el backend}

El script \archivo{ocr/LecturaTickets\_v2.py} implementa un pipeline de
6 etapas ejecutado como subproceso de Node.js:

\begin{enumerate}
  \item \textbf{Preprocesamiento} (OpenCV): upscaling a mínimo 1\,200 px,
        corrección de iluminación, deskew, binarización adaptativa
        gaussiana, denoising y sharpening.
  \item \textbf{OCR con Tesseract}: modo LSTM (\campo{oem 3}), bloque de
        texto (\campo{psm 6}), idiomas \campo{spa+eng}. Retorna texto y
        confianza promedio por palabra.
  \item \textbf{Evaluación de calidad}: dos compuertas independientes.
  \item \textbf{Respaldo Vision + Claude} (si aplica).
  \item \textbf{Parsing por regex}: extracción de productos, total y fecha.
  \item \textbf{Categorización}: pipeline de 3 pasos por producto.
\end{enumerate}

\subsubsection*{Compuertas de calidad para activar el respaldo}

\begin{table}[H]
  \centering
  \caption{Compuertas de calidad del pipeline OCR}
  \label{tab:compuertas-ocr}
  \begin{tabular}{llp{6cm}}
    \toprule
    \textbf{Compuerta} & \textbf{Condición} & \textbf{Descripción} \\
    \midrule
    1 — Técnica &
      \campo{conf < 0.72} O \campo{palabras < 15} &
      Baja confianza por palabra o texto escaso. Se activa antes del
      parsing. \\[4pt]
    2 — Estructural &
      Sin total Y sin productos &
      Tesseract pasó la Compuerta 1 pero el parsing no encontró datos
      útiles. Se activa después del parsing. \\
    \bottomrule
  \end{tabular}
\end{table}

\subsubsection*{Respaldo con Google Cloud Vision + Claude}

Cuando se activa cualquier compuerta, la función
\campo{\_enhance\_extraction()}:

\begin{enumerate}
  \item Codifica la imagen en base64 y la envía a la API de Google Cloud
        Vision
  \item Si se obtiene texto, lo envía a Claude Haiku con un prompt que
        estructura el ticket en JSON
        (\campo{store\_name, date, total, products}).
  \item El JSON enriquecido se deserializa al formato interno del pipeline.
  \item La confianza se establece en 0.95 si Vision fue exitoso.
\end{enumerate}

Si las claves API no están disponibles o cualquier paso falla, la función
retorna el texto original de Tesseract sin lanzar excepción.

\subsubsection*{Categorización automática de productos}

Cada producto detectado pasa por un pipeline de 3 pasos:

\begin{enumerate}
  \item \textbf{categorias.json}: comparación del nombre del producto
        contra palabras clave organizadas por categoría. Match por
        substring exacto o similitud \campo{SequenceMatcher}
        $\geq 0.75$.
  \item \textbf{Historial del usuario}: si el Paso~1 no encontró
        categoría, se consultan los productos previos del mismo usuario
        en \campo{Productos\_Tickets}. Match por substring o similitud
        $\geq 0.70$.
  \item \textbf{Fallback}: si ninguno encontró coincidencia, el producto
        se asigna a la categoría ``Otros'' con
        \campo{id\_categoria = null}.
\end{enumerate}

\begin{lstlisting}[style=python,
    caption={Pipeline de categorización de productos — Paso 1}]
# Pipeline de categorizacion -- Paso 1 (categorias.json)
def _classify_by_json(name, categories):
    name_up = name.upper()
    for cat, keywords in categories.items():
        if cat == "OTROS": continue
        for kw in keywords:
            if kw.upper() in name_up:
                return cat          # Coincidencia exacta por substring
            if SequenceMatcher(None, name_up, kw.upper()).ratio() >= 0.75:
                return cat          # Coincidencia fuzzy
    return None                     # Sin coincidencia
\end{lstlisting}

\subsubsection*{Métrica de confianza del ticket}

La función \campo{calculate\_confidence()} combina la confianza técnica
de Tesseract con señales estructurales del parsing. El valor final se
normaliza al rango $[0.0,\ 1.0]$ con \campo{max(0, min(score, 1))}.

\begin{longtable}{lrp{6.5cm}}
  \caption{Componentes de la métrica de confianza OCR}
  \label{tab:confianza-ocr} \\
  \toprule
  \textbf{Componente} & \textbf{Peso} & \textbf{Justificación} \\
  \midrule
  \endfirsthead
  \toprule
  \textbf{Componente} & \textbf{Peso} & \textbf{Justificación} \\
  \midrule
  \endhead
  \campo{ocr\_conf × 0.30}  & 30\% &
    Confianza promedio por palabra de Tesseract (normalizada a 0--1). \\[3pt]
  \campo{total\_found}       & +25\% / $-$15\% &
    El total es el campo crítico de todo ticket. \\[3pt]
  \campo{n\_products > 0}    & +20\% / $-$15\% &
    Los productos son el contenido principal. \\[3pt]
  \campo{store\_found}       & +10\% &
    Contextual; confirma que el documento es un ticket. \\[3pt]
  \campo{date\_found}        & +10\% / $-$5\% &
    Permite clasificación temporal del gasto. \\[3pt]
  \campo{word\_count < 10}   & $-$10\% &
    Imagen casi ilegible aunque las pocas palabras sean nítidas. \\[3pt]
  \campo{word\_count < 25}   & $-$5\% &
    Imagen con texto escaso (posible ticket parcial). \\
  \bottomrule
\end{longtable}

\subsubsection{Servicio de tickets}

\begin{table}[H]
  \centering
  \caption{Métodos del servicio de tickets}
  \label{tab:servicio-tickets}
  \begin{tabular}{lp{8.5cm}}
    \toprule
    \textbf{Método} & \textbf{Descripción} \\
    \midrule
    \campo{scanTicket(imagePaths)} &
      POST multipart a \endpoint{/tickets/scan}. Timeout 60 s. \\
    \campo{getAll()} &
      GET \endpoint{/tickets}. Lista de tickets del usuario. \\
    \campo{getById(id)} &
      GET \endpoint{/tickets/:id}. Ticket con imagen en base64. \\
    \campo{getProducts(ticketId)} &
      GET \endpoint{/tickets/:id/products}. Productos con categoría. \\
    \campo{delete(id)} &
      DELETE \endpoint{/tickets/:id}. Elimina ticket y productos. \\
    \bottomrule
  \end{tabular}
\end{table}

\subsubsection{Consideraciones técnicas}

\begin{itemize}
  \item \textbf{Idempotencia}: el backend genera un hash SHA-256 de las
        imágenes para evitar procesar el mismo ticket dos veces (error
        409 en duplicados).
  \item \textbf{Limpieza}: los archivos temporales de imagen se eliminan
        del servidor tras procesar el ticket.
  \item \textbf{Transacciones DB}: el ticket y sus productos se insertan
        en una transacción con \campo{rollback} ante cualquier falla.
  \item \textbf{Manejo de errores en Flutter}: si el OCR lanza excepción,
        la pantalla de procesamiento permite reintentar sin necesidad de
        fotografiar nuevamente el ticket.
\end{itemize}


% ──────────────────────────────────────────────────────────────
\subsection{Módulo de categorías}
% ──────────────────────────────────────────────────────────────

\subsubsection{Descripción general}

El módulo de categorías provee la estructura de clasificación que sustenta
todos los demás módulos. Cada gasto, producto de ticket, presupuesto e
ingreso queda asociado a una categoría del usuario. El sistema distingue
dos tipos: \textbf{predefinidas} (creadas durante el onboarding) y
\textbf{personalizadas} (creadas libremente por el usuario).

Las categorías predefinidas por defecto son: Alimentos, Entretenimiento,
Transporte, Ahorro y Ropa.

\subsubsection{Modelo de datos}

\begin{lstlisting}[style=dart,
    caption={Modelo de datos Category en Flutter}]
class Category {
  final int    id;
  final int    idUsuario;
  final String nombreCategoria;
  final String tipoCategoria;  // 'predefinida' | 'personalizada'
  final String? icono;         // Nombre del icono (ej. 'restaurant')
  final String? color;         // Color hex (ej. '#13853d')

  IconData getIconData();      // Mapea icono a FlutterIcon
  Color    getColor();         // Parsea hex a Color de Flutter
}
\end{lstlisting}

El método \campo{getColor()} parsea cadenas hexadecimales en formato
\campo{\#RRGGBB} o \campo{0xFFRRGGBB}. Si el formato no es reconocido,
retorna un color neutro por defecto (\campo{Colors.blueGrey}).

\subsubsection{Pantallas}

\subsubsection*{CategoriesScreen — Listado}

Muestra las categorías en dos secciones diferenciadas:

\begin{itemize}
  \item \textbf{Predefinidas}: grid de 3 columnas con icono y nombre.
        No son editables ni eliminables.
  \item \textbf{Personalizadas}: lista con icono, nombre y botón de
        edición. El usuario puede crear, editar y eliminar las suyas.
\end{itemize}

\subsubsection*{NewCategoryScreen — Creación}

Formulario con tres campos:

\begin{enumerate}
  \item Nombre de la categoría (obligatorio, máx. 50 caracteres).
  \item Icono: selector de cuadrícula con 20+ opciones predefinidas.
  \item Color: selector de cuadrícula con 12 colores disponibles.
\end{enumerate}

\subsubsection*{EditCategoryScreen — Edición}

Igual que la creación, pero precarga los valores actuales de la
categoría recibida como argumento de navegación.

\subsubsection{Servicio de categorías}

\begin{table}[H]
  \centering
  \caption{Métodos del servicio de categorías}
  \label{tab:servicio-categorias}
  \begin{tabular}{lp{8.5cm}}
    \toprule
    \textbf{Método} & \textbf{Descripción} \\
    \midrule
    \campo{getAll()} &
      GET \endpoint{/categories}. Todas las categorías del usuario. \\
    \campo{getById(id)} &
      GET \endpoint{/categories/:id}. \\
    \campo{create(\{nombre, icono, color\})} &
      POST \endpoint{/categories}. Crea categoría personalizada. \\
    \campo{update(\{id, nombre, icono, color\})} &
      PUT \endpoint{/categories/:id}. \\
    \campo{delete(id)} &
      DELETE \endpoint{/categories/:id}. \\
    \campo{getPredefined()} &
      Filtra localmente \campo{tipoCategoria == 'predefinida'}. \\
    \campo{getCustom()} &
      Filtra localmente \campo{tipoCategoria == 'personalizada'}. \\
    \bottomrule
  \end{tabular}
\end{table}

\subsubsection{Flujo funcional}

\begin{enumerate}
  \item El usuario navega a la pantalla de Categorías desde el módulo de
        Presupuestos o desde el menú de configuración.
  \item Ve el listado separado en dos secciones.
  \item Para crear: presiona el botón ``+'' y completa el formulario.
  \item El servicio llama al endpoint POST y recarga la lista en caso de
        éxito.
  \item La nueva categoría queda disponible inmediatamente para asignar
        en presupuestos, gastos y productos de tickets.
\end{enumerate}

\subsubsection{Consideraciones técnicas}

\begin{itemize}
  \item Las categorías predefinidas no pueden eliminarse desde el
        frontend (el botón de edición no aparece en la sección
        predefinidas).
  \item En caso de error al cargar, se muestra un estado de error con
        botón de reintentar en lugar de pantalla vacía.
  \item El servicio retorna lista vacía (no lanza excepción) ante
        cualquier error HTTP, para no bloquear formularios que dependan
        de las categorías.
\end{itemize}


% ──────────────────────────────────────────────────────────────
\subsection{Módulo de presupuestos}
% ──────────────────────────────────────────────────────────────

\subsubsection{Descripción general}

El módulo de presupuestos permite al usuario asignar un límite de gasto
mensual por categoría. El sistema calcula automáticamente cuánto se ha
gastado en cada categoría durante el mes en curso y emite alertas cuando
el presupuesto está cerca de agotarse o ha sido superado.

Cada presupuesto cubre el rango de fechas del mes calendario actual
(del día~1 al último día del mes), calculado automáticamente tanto en el
frontend como en el backend.

\subsubsection{Modelo de datos}

\begin{lstlisting}[style=dart,
    caption={Modelo de datos Budget en Flutter}]
class Budget {
  final int      id;
  final int      idCategoria;
  final double   montoAsignado;
  final double   montoUtilizado;   // Actualizado por el backend
  final String   periodo;          // 'mensual' (fijo)
  final DateTime fechaInicio;
  final DateTime fechaFin;
  final bool     activo;

  // Propiedades calculadas
  double get porcentajeUtilizado =>
      (montoUtilizado / montoAsignado * 100).clamp(0, double.infinity);

  double get montoDisponible =>
      (montoAsignado - montoUtilizado).clamp(0, double.infinity);
}
\end{lstlisting}

\subsubsection{Pantallas}

\subsubsection*{BudgetsScreen — Panel de presupuestos}

La pantalla principal del módulo presenta:

\begin{itemize}
  \item \textbf{Tarjeta resumen global}: presupuesto total asignado,
        total gastado y porcentaje de uso del mes.
  \item \textbf{Barra de progreso global}: cambia a rojo cuando el uso
        supera el 90\%.
  \item \textbf{Tarjetas por presupuesto}: icono de categoría, nombre,
        barra de progreso individual, monto gastado vs.\ asignado,
        porcentaje y botón de edición.
  \item \textbf{Estado vacío}: con ícono y llamada a la acción para
        crear el primer presupuesto.
\end{itemize}

\begin{lstlisting}[style=dart,
    caption={Cálculo de totales en BudgetsScreen}]
// Calculo de totales en la pantalla
double get _totalBudget =>
    _budgets.fold(0, (s, b) => s + b.montoAsignado);

double get _totalSpent =>
    _budgets.fold(0, (s, b) => s + b.montoUtilizado);

double get _percentageUsed =>
    _totalBudget > 0 ? (_totalSpent / _totalBudget * 100) : 0;
\end{lstlisting}

\subsubsection*{NewBudgetScreen — Creación de presupuesto}

El formulario presenta un grid de categorías disponibles para seleccionar
(2 columnas, con resaltado visual de la seleccionada) y un campo de monto.

Validaciones implementadas:

\begin{enumerate}
  \item Monto requerido y mayor a cero.
  \item Categoría seleccionada obligatoria.
  \item Verificación de presupuesto duplicado: si la categoría ya tiene
        un presupuesto activo, se muestra un aviso y no se crea uno
        nuevo.
\end{enumerate}

\begin{lstlisting}[style=dart,
    caption={Verificación de presupuesto duplicado}]
// Verificacion de presupuesto duplicado antes de crear
final existingBudget = await _budgetsService.getByCategory(
    _selectedCategoryId!);
if (existingBudget != null) {
  ScaffoldMessenger.of(context).showSnackBar(
    const SnackBar(content: Text(
        'Esta categoria ya tiene un presupuesto asignado')));
  return;
}
\end{lstlisting}

\subsubsection*{EditBudgetScreen — Edición}

Permite modificar el monto asignado, la categoría y el estado
activo/inactivo de un presupuesto existente. Recibe el \campo{budgetId}
como argumento de navegación y carga los datos actuales del presupuesto.

\subsubsection{Servicio de presupuestos}

\begin{table}[H]
  \centering
  \caption{Métodos del servicio de presupuestos}
  \label{tab:servicio-presupuestos}
  \begin{tabular}{lp{8.5cm}}
    \toprule
    \textbf{Método} & \textbf{Descripción} \\
    \midrule
    \campo{getAll()} &
      GET \endpoint{/budgets}. Incluye \campo{porcentaje\_usado}
      precalculado. \\
    \campo{getByCategory(categoryId)} &
      GET \endpoint{/budgets/category/:id}. Para validar duplicados. \\
    \campo{create(data)} &
      POST \endpoint{/budgets}. Fechas auto-calculadas por backend. \\
    \campo{update(id, data)} &
      PUT \endpoint{/budgets/:id}. \\
    \campo{delete(id)} &
      DELETE \endpoint{/budgets/:id}. \\
    \campo{getActiveBudgets()} &
      Filtra localmente \campo{activo == true}. \\
    \campo{getBudgetProgress(budget)} &
      Calcula estado: safe / warning / danger / exceeded. \\
    \bottomrule
  \end{tabular}
\end{table}

\subsubsection{Estados de progreso de presupuesto}

\begin{table}[H]
  \centering
  \caption{Estados de progreso de un presupuesto}
  \label{tab:estados-presupuesto}
  \begin{tabular}{llp{6cm}}
    \toprule
    \textbf{Estado} & \textbf{Umbral} & \textbf{Indicación visual} \\
    \midrule
    \campo{safe}     & $< 50\%$  & Barra azul \\
    \campo{warning}  & 50--80\%  & Barra naranja \\
    \campo{danger}   & 80--100\% & Barra roja \\
    \campo{exceeded} & $> 100\%$ & Barra roja + alerta \\
    \bottomrule
  \end{tabular}
\end{table}

\subsubsection{Flujo funcional}

\begin{enumerate}
  \item El usuario accede al módulo de Presupuestos desde el
        \campo{BottomAppBar} en la pantalla principal.
  \item La pantalla carga todos los presupuestos activos del mes.
  \item El usuario puede ver el progreso individual de cada categoría.
  \item Si desea crear uno nuevo, selecciona la categoría y el monto
        límite.
  \item El backend auto-asigna las fechas del mes actual.
  \item Cuando un gasto nuevo se registra (OCR o manual), el backend
        actualiza el campo \campo{monto\_utilizado} del presupuesto
        correspondiente.
  \item Al rebasar el 80\% o el 100\%, el sistema de alertas genera una
        notificación para el usuario.
\end{enumerate}

\subsubsection{Consideraciones técnicas}

\begin{itemize}
  \item El campo \campo{montoUtilizado} es calculado y actualizado por
        el backend; el frontend solo lo lee.
  \item El mapa \campo{Map<int, Category>} en la pantalla permite
        búsquedas en O(1) sin recorrer la lista completa de categorías
        en cada render.
  \item Al regresar de la creación/edición con \campo{result == true},
        la pantalla recarga automáticamente la lista completa.
  \item El módulo llama a
        \campo{NotificationTriggers().checkAlertsOnce()} al cargar para
        verificar si hay alertas pendientes de presupuesto.
\end{itemize}


% ──────────────────────────────────────────────────────────────
\subsection{Módulo de deudas}
% ──────────────────────────────────────────────────────────────

\subsubsection{Descripción general}

El módulo de deudas gestiona créditos e instrumentos de deuda con soporte
para la \textbf{fórmula de amortización francesa}. Permite registrar el
capital inicial, la tasa de interés anual, el plazo en meses y la fecha
del próximo pago. El backend calcula el pago mensual, la tasa mensual
equivalente y el saldo insoluto tras cada abono. El usuario puede realizar
pagos ordinarios y abonos extraordinarios a capital.

\subsubsection{Modelo de datos}

\begin{lstlisting}[style=dart,
    caption={Modelo de datos Debt en Flutter}]
class Debt {
  final int      id;
  final String   nombreDeuda;
  final double   capitalInicial;
  final double   tasaAnual;
  final double   tasaMensual;      // tasaAnual / 12
  final double   saldoInsoluto;    // Saldo pendiente post-pagos
  final double   pagoMensual;      // Calculado por backend
  final double   montoPagado;      // Acumulado de pagos
  final int      duracionMeses;
  final DateTime fechaSiguientePago;
  final bool     activa;

  // Propiedades calculadas en cliente
  double get porcentajePagado  => montoPagado / capitalInicial * 100;
  double get montoPendiente    => (capitalInicial - montoPagado)
                                    .clamp(0, double.infinity);
  double get interesTotal      => (pagoMensual * duracionMeses)
                                    - capitalInicial;
  bool   get tieneInteres      => tasaAnual > 0;
  int    get diasHastaPago     => fechaSiguientePago
                                    .difference(DateTime.now()).inDays;
}
\end{lstlisting}

\subsubsection{Fórmula de amortización francesa}

El pago mensual se calcula usando la fórmula estándar de cuota fija:

\[
  M = C \cdot \frac{r}{1 - (1 + r)^{-n}}
\]

donde $C$ es el capital inicial, $r$ es la tasa mensual
($r = \text{tasa\_anual} / 1200$) y $n$ es el número de mensualidades.
El frontend calcula esta fórmula en tiempo real para mostrar una vista
previa antes de guardar la deuda:

\begin{lstlisting}[style=dart,
    caption={Cálculo del pago mensual con amortización francesa}]
// Calculo del pago mensual con amortizacion francesa
double get _tasaMensual => _tasaAnual / 1200;

double get _pagoMensual {
  if (_capital <= 0 || _meses <= 0) return 0;
  if (_tasaMensual == 0) return _capital / _meses;
  final r = _tasaMensual;
  final n = _meses;
  return _capital * r / (1 - pow(1 + r, -n));
}
\end{lstlisting}

Para deudas sin interés (\campo{tasaAnual == 0}), el cálculo simplifica
a la división directa del capital entre el número de meses.

\subsubsection{Pantallas}

\subsubsection*{DebtsScreen — Panel de deudas}

Muestra el estado consolidado de todas las deudas del usuario:

\begin{itemize}
  \item Tarjeta resumen: deuda total (naranja), total pagado (verde) y
        porcentaje pagado global.
  \item Por cada deuda: nombre, barra de progreso (rojo $<$ 50\%,
        naranja $\geq$ 50\%), monto pagado vs.\ total, días al próximo
        pago y botones de pago y edición.
  \item El botón de pago abre un diálogo modal donde el usuario ingresa
        el monto a pagar.
\end{itemize}

\subsubsection*{NewDebtScreen — Creación}

El formulario captura los datos de la deuda y actualiza en tiempo real
una tarjeta de resumen financiero:

\begin{itemize}
  \item Nombre de la deuda.
  \item Capital inicial (monto del crédito).
  \item Plazo en meses (campo de texto libre, por defecto 12).
  \item Tasa de interés anual en porcentaje (por defecto 0\%).
  \item Fecha del próximo pago (DatePicker, por defecto +30 días).
\end{itemize}

La tarjeta de resumen en tiempo real muestra: tasa mensual equivalente,
pago mensual calculado e interés total del crédito.

\subsubsection*{EditDebtScreen — Edición y abono a capital}

Además de editar los campos del formulario, incluye un botón exclusivo
para registrar \textbf{pagos anticipados a capital}. Este tipo de pago
reduce el saldo insoluto directamente sin afectar la cuota de intereses
del período, lo que disminuye las mensualidades restantes.

\subsubsection{Servicio de deudas}

\begin{table}[H]
  \centering
  \caption{Métodos del servicio de deudas}
  \label{tab:servicio-deudas}
  \begin{tabular}{lp{8.5cm}}
    \toprule
    \textbf{Método} & \textbf{Descripción} \\
    \midrule
    \campo{getAll()} &
      GET \endpoint{/debts}. Todas las deudas del usuario. \\
    \campo{create(data)} &
      POST \endpoint{/debts}. El backend calcula \campo{pago\_mensual}
      y \campo{tasa\_mensual}. \\
    \campo{update(id, data)} &
      PUT \endpoint{/debts/:id}. \\
    \campo{registerPayment(id, monto)} &
      POST \endpoint{/debts/:id/pay}. Registra abono ordinario.
      Retorna deuda con \campo{saldo\_insoluto} actualizado. \\
    \campo{registerCapitalPayment(id, monto)} &
      POST \endpoint{/debts/:id/capital-payment}. Abono
      extraordinario a capital. \\
    \campo{delete(id)} &
      DELETE \endpoint{/debts/:id}. \\
    \campo{getActiveDebts()} &
      Filtra localmente \campo{activa == true}. \\
    \campo{getUpcomingPayments(\{days\})} &
      Filtra deudas con \campo{diasHastaPago $\leq$ days}. \\
    \campo{getDebtProgress(debt)} &
      Retorna estado: critical / warning / progress / almost / paid. \\
    \bottomrule
  \end{tabular}
\end{table}

\subsubsection{Estados de progreso de deuda}

\begin{table}[H]
  \centering
  \caption{Estados de progreso de una deuda}
  \label{tab:estados-deuda}
  \begin{tabular}{llp{6cm}}
    \toprule
    \textbf{Estado} & \textbf{Umbral} & \textbf{Indicación visual} \\
    \midrule
    \campo{critical} & $< 25\%$ pagado    & Rojo \\
    \campo{warning}  & 25--50\% pagado    & Naranja \\
    \campo{progress} & 50--75\% pagado    & Amarillo \\
    \campo{almost}   & 75--100\% pagado   & Verde claro \\
    \campo{paid}     & $= 100\%$          & Verde \\
    \bottomrule
  \end{tabular}
\end{table}

\subsubsection{Flujo funcional}

\begin{enumerate}
  \item El usuario accede al módulo de Deudas desde el
        \campo{BottomAppBar}.
  \item El panel carga y muestra todas las deudas activas con su estado.
  \item Para registrar un pago: botón verde en la tarjeta de deuda
        $\to$ diálogo modal con campo de monto $\to$ llamada al endpoint
        \endpoint{/debts/:id/pay} $\to$ recarga de la lista.
  \item Para abono a capital: editar deuda $\to$ botón de abono a
        capital $\to$ diálogo con monto $\to$ llamada a
        \endpoint{/debts/:id/capital-payment}.
  \item El módulo verifica alertas de pagos próximos al cargar
        (\campo{NotificationTriggers}).
\end{enumerate}

\subsubsection{Consideraciones técnicas}

\begin{itemize}
  \item \textbf{Cálculo en frontend vs.\ backend}: la fórmula de
        amortización se ejecuta en tiempo real en el cliente para la
        vista previa, pero el backend recalcula y persiste los valores
        canónicos al guardar.
  \item \textbf{Validaciones}: el monto de pago debe ser mayor que cero
        y no puede exceder el saldo pendiente (\campo{montoPendiente}).
  \item \textbf{Deuda sin interés}: la fórmula degenera a división
        simple; la UI oculta la línea de ``Interés total'' cuando la
        tasa es cero.
  \item \textbf{Plazo variable}: el campo de plazo es un
        \campo{TextField} libre para facilitar el ingreso de valores
        grandes.
  \item Los días hasta el próximo pago se calculan en el cliente con
        \campo{DateTime.difference} para evitar tráfico innecesario.
\end{itemize}


% ══════════════════════════════════════════════════════════════
\section{Aplicación web}
% ══════════════════════════════════════════════════════════════


\subsection{Arquitectura general de la aplicación web}

\subsubsection{Stack tecnológico}

\begin{longtable}{>{\bfseries\small}p{3.5cm} p{10cm}}
  \caption{Stack tecnológico de la aplicación web FINARA}
  \label{tab:stack-web} \\
  \toprule
  Componente & Descripción \\
  \midrule
  \endfirsthead
  \toprule
  Componente & Descripción \\
  \midrule
  \endhead
  Next.js + React & Framework principal del frontend web (App Router), con render híbrido y estructura modular por rutas en \archivo{src/app}. \\[4pt]
  TypeScript & Tipado estático para componentes, servicios y modelos de datos de la capa web. \\[4pt]
  Tailwind CSS & Sistema de utilidades para estilos y layout responsivo en paneles, tarjetas, tablas y modales. \\[4pt]
  next-auth & Manejo de sesión/autenticación en frontend mediante \campo{SessionProvider} y \campo{useSession}. \\[4pt]
  jsPDF + autotable & Generación local de reportes ejecutivos en PDF desde la vista de dashboard. \\[4pt]
  xlsx & Generación local de reportes en Excel (.xlsx) con hojas de resumen y detalle. \\[4pt]
  Node.js + Express + MySQL & Backend REST consumido por la web para operaciones transaccionales, recomendaciones y persistencia. \\[4pt]
  OCR backend (Python/Tesseract + Vision/Claude) & Pipeline de extracción consumido por endpoints de tickets; la web integra el flujo de escaneo mediante API REST. \\[4pt]
  Cliente HTTP centralizado & \archivo{src/lib/backendApi.ts}: inyección de token, envelope \campo{\{success, data\}}, control de errores y redirección por sesión expirada. \\
  \bottomrule
\end{longtable}

\subsubsection{Estructura del proyecto web}

La aplicación web se organiza por rutas y módulos de dominio dentro de
\archivo{src/app/}, con componentes reutilizables y servicios desacoplados en
\archivo{src/lib/}.

\begin{lstlisting}[style=codigobase, language={},
    caption={Estructura de directorios relevante en la aplicación web}]
src/
|-- app/
|   |-- dashboard/
|   |   |-- page.tsx                  (dashboard de analisis)
|   |   |-- recomendaciones/page.tsx  (modulo recomendaciones)
|   |   |-- reportes/page.tsx         (historial de reportes)
|   |   `-- transacciones/page.tsx    (registro/consulta de movimientos)
|   |-- api/
|   |   `-- scan-ticket/route.ts      (endpoint OCR interno en Next)
|   `-- components/
|       |-- FloatingActionButton.tsx  (alta de ticket/ingresos)
|       |-- TicketPreview.tsx         (revision de OCR)
|       `-- ToastStack.tsx            (feedback de eventos)
`-- lib/
    |-- backendApi.ts                 (cliente HTTP centralizado)
    |-- scanTicketService.ts          (servicio escaneo ticket web)
    `-- reports/
        |-- reportApi.ts              (API historial reportes)
        |-- generators/pdfGenerator.ts
        |-- generators/excelGenerator.ts
        `-- utils/analysisUtils.ts
\end{lstlisting}

\subsubsection{Comunicación con el backend}

Toda la integración de la web con backend se realiza vía \textbf{REST API}.
El helper \campo{backendRequest()} concentra comportamiento común:

\begin{itemize}
  \item Inyección automática de \campo{Authorization: Bearer <token>}.
  \item Parseo de envelopes en formato \campo{\{success, data, error\}}.
  \item Manejo de errores HTTP y errores de negocio.
  \item Redirección automática a login ante expiración de sesión.
  \item Normalización de URL base con \campo{getBackendBaseUrl()}.
\end{itemize}

\subsubsection{Manejo de estado}

La web utiliza estado local con React:
\begin{itemize}
  \item \campo{useState}, \campo{useEffect}, \campo{useMemo} para estado de pantalla y derivados.
  \item Estados de carga/error por módulo (dashboard, recomendaciones, reportes, escaneo).
  \item Estado de sesión con \campo{next-auth} y proveedor \campo{SessionWrapper}.
\end{itemize}

No se usa una librería global dedicada; el estado se modela
por módulo y servicios.


% ──────────────────────────────────────────────────────────────
\subsection{Módulo de registro de tickets digitales (flujo web)}
% ──────────────────────────────────────────────────────────────

\subsubsection{Descripción general}

Este módulo permite registrar tickets desde la web subiendo un archivo
\textbf{PDF o imagen} (JPG/JPEG/PNG), procesarlo mediante backend OCR,
revisar los datos detectados y confirmar el guardado final del ticket.

\begin{notabox}
  A diferencia del flujo móvil, en la web no existe etapa de captura con
  cámara nativa El flujo inicia con selección de
  archivo desde el sistema operativo.
\end{notabox}

\subsubsection*{Flujo completo de operación (web)}

\begin{enumerate}
  \item \textbf{Selección de archivo} (FAB): el usuario elige PDF/JPG/PNG.
  \item \textbf{Escaneo OCR}: frontend envía \campo{FormData} al endpoint
        \endpoint{/tickets/scan-web} usando \campo{scanTicketByWebFile()}.
  \item \textbf{Previsualización}: \campo{TicketPreview} muestra imagen/PDF,
        establecimiento, fecha, total, ítems, confianza OCR y método de extracción.
  \item \textbf{Confirmación y guardado}: se envía payload final al endpoint
        de persistencia (por metadata de escaneo o \endpoint{/tickets/manual}).
\end{enumerate}

\subsubsection{Validaciones implementadas}

\begin{itemize}
  \item Tipos permitidos: \campo{application/pdf}, \campo{image/jpeg},
        \campo{image/jpg}, \campo{image/png}.
  \item Sesión/token obligatorios para escanear y guardar.
  \item Verificación de categorías disponibles antes de confirmar el ticket.
  \item Validación de productos con categoría válida al guardar.
\end{itemize}

\subsubsection{Componentes y servicios principales}

\begin{table}[H]
  \centering
  \caption{Elementos principales del módulo de ticket digital web}
  \label{tab:modulo-ticket-web}
  \begin{tabular}{lp{8.5cm}}
    \toprule
    \textbf{Componente/Servicio} & \textbf{Responsabilidad} \\
    \midrule
    \archivo{FloatingActionButton.tsx} &
      Orquesta selección de archivo, invoca escaneo, gestiona edición/confirmación y ejecuta guardado final. \\
    \archivo{TicketPreview.tsx} &
      Renderiza previsualización del ticket (iframe para PDF o imagen), estado OCR y detalle de productos detectados. \\
    \archivo{scanTicketService.ts} &
      Implementa \campo{scanTicketByWebFile()} y normaliza payload OCR del backend. \\
    \archivo{backendApi.ts} &
      Cliente HTTP para operaciones autenticadas y manejo de errores comunes. \\
    \bottomrule
  \end{tabular}
\end{table}

\subsubsection{Persistencia y payload de guardado}

Cuando el ticket proviene de escaneo, el frontend construye payload con:
\begin{itemize}
  \item datos del ticket (tienda, fecha, total),
  \item productos editables (nombre, precio, categoría),
  \item metadatos OCR (\campo{confianza\_ocr}, \campo{request\_id}),
  \item imagen del ticket en base64 y MIME cuando está disponible.
\end{itemize}

\begin{lstlisting}[style=codigobase,
    caption={Guardado final de ticket escaneado (flujo web)}]
const payloadBody = {
  nombre_tienda: establecimiento,
  total: Number(calcularTotal()) || 0,
  fecha_compra: apiDate,
  id_categoria: items[0]?.id_categoria,
  productos: items.map(item => ({
    nombre_producto: item.nombre_producto,
    precio: item.precio,
    id_categoria: item.id_categoria,
  })),
  request_id: requestId || undefined,
  confianza_ocr: scanConfidenceOcr ?? undefined,
  imagen_ticket_base64: extractedImage.base64 || undefined,
  imagen_ticket_mime: extractedImage.mime || undefined,
};
\end{lstlisting}

\
% ──────────────────────────────────────────────────────────────
\subsection{Módulo dashboard inicial (análisis y gráficas)}
% ──────────────────────────────────────────────────────────────

\subsubsection{Descripción general}

El dashboard web es una vista de análisis financiero para comprensión rápida
por período. Integra KPIs, tendencias y distribución por categoría
para apoyar lectura analítica (más orientado a análisis).

\subsubsection{Objetivo}

Concentrar indicadores financieros del usuario en una sola vista para detectar
comportamientos de gasto/ingreso, balance y uso de presupuesto.

\subsubsection{Carga de datos y agregaciones}

La función \campo{loadDashboardData()} consume en paralelo:
\begin{itemize}
  \item \endpoint{/dashboard/summary}
  \item \endpoint{/incomes}
  \item \endpoint{/expenses}
  \item \endpoint{/debts}
  \item \endpoint{/budgets}
  \item \endpoint{/categories}
\end{itemize}

Con esos datos se calculan:
\begin{itemize}
  \item ingresos y gastos del período,
  \item deuda activa y deuda global,
  \item porcentaje de presupuesto utilizado,
  \item series semanales (mes) y mensuales (año),
  \item agregación de gastos por categoría.
\end{itemize}

\subsubsection{Gráficas implementadas}

\begin{table}[H]
  \centering
  \caption{Gráficas del dashboard web y su propósito analítico}
  \label{tab:graf-dashboard-web}
  \begin{tabular}{lp{8.5cm}}
    \toprule
    \textbf{Gráfica} & \textbf{Uso analítico} \\
    \midrule
    \campo{BarChart Ingresos vs Gastos} &
      Comparar comportamiento financiero por semana (modo mes) o por mes (modo año). \\
    \campo{Flujo de ingresos} &
      Visualizar el flujo del dinero del usuario mes con mes. \\
    \campo{BarChart Gastos por categoría} &
      Identificar categorías con mayor concentración de egresos. \\
    \campo{PieChart de distribución} &
      Leer balance general y estado de presupuesto de forma rápida. \\
    \bottomrule
  \end{tabular}
\end{table}

\subsubsection{Interacción funcional}

\begin{enumerate}
  \item El usuario selecciona mes y año.
  \item Aplica filtros y recarga la data del período.
  \item Puede cambiar entre vista mensual y anual.
  \item Puede disparar generación de PDF/Excel con el período activo.
\end{enumerate}

\subsubsection{Consideraciones técnicas}

\begin{itemize}
  \item Uso de \campo{Promise.all} para mejorar tiempo de carga inicial.
  \item Uso de \campo{useMemo} para derivados de series y categoría.
  \item Manejo explícito de estados \campo{loadingData} y \campo{loadError}.
  \item Revalidación por eventos de actualización de ingresos
        (\campo{finara:income-updated}).
\end{itemize}


% ──────────────────────────────────────────────────────────────
\subsection{Módulo de generación de reportes}
% ──────────────────────────────────────────────────────────────

\subsubsection{Descripción general}

La aplicación web se manejan reportes:
  \textbf{Generación local} desde dashboard (descarga inmediata PDF/Excel).
  


\subsubsection{Generación de PDF y Excel desde dashboard}

Desde \archivo{src/app/dashboard/page.tsx}:
\begin{itemize}
  \item \campo{handleDownloadPDF()} invoca
        \campo{generateExecutiveSummaryPDF(...)}.
  \item \campo{handleDownloadExcel()} invoca
        \campo{generateExecutiveSummaryExcel(...)}.
\end{itemize}

Ambos incluyen contexto del período filtrado y recomendaciones del endpoint
\endpoint{/recommendations} para enriquecer resumen ejecutivo.




\subsubsection{Consideraciones técnicas}

\begin{itemize}
  \item En dashboard se valida que existan movimientos en el período antes de generar archivo.

\end{itemize}


% ──────────────────────────────────────────────────────────────
\subsection{Módulo de recomendaciones}
% ──────────────────────────────────────────────────────────────

\subsubsection{Descripción general}

Este módulo muestra recomendaciones financieras del backend para un período
(mes/año), integrando riesgos, recomendaciones y acciones sugeridas.

\subsubsection{Objetivo}

Traducir el comportamiento financiero del usuario en recomendaciones accionables
para mejor control de gastos y balance mensual.

\subsubsection{Pantalla principal y flujo}

Implementado en \archivo{src/app/dashboard/recomendaciones/page.tsx}:

\begin{enumerate}
  \item El sistema toma el periodo actual.
  \item La pantalla consulta en paralelo \endpoint{/recommendations},
        \endpoint{/incomes} y \endpoint{/expenses}.
  \item Se calculan métricas del período (ingresos fijos, ingresos únicos,
        ingresos totales, gastos y balance).
  \item Se renderizan tarjetas de resumen, riesgos y acciones priorizadas.
\end{enumerate}

\subsubsection{Estructura de datos mostrada}

\begin{itemize}
  \item \campo{resumen}
  \item \campo{riesgos[]}
  \item \campo{recomendaciones[]}
  \item \campo{acciones\_sugeridas[]}
  \item \campo{nota}
\end{itemize}

\subsubsection{Integración con otros módulos}

El módulo de recomendaciones no solo se consulta en su pantalla propia;
también se reutiliza en el dashboard para enriquecer contenido de reportes
PDF/Excel (resumen, riesgos y sugerencias).

\subsubsection{Consideraciones técnicas}

\begin{itemize}
  \item Normalización defensiva del payload para soportar variantes
        \campo{\{data: ...\}} o payload directo.
  \item Redirección a login ante errores de autenticación (401/JWT).
  \item Estados diferenciados de \campo{loading}, \campo{refreshing} y \campo{error}.
\end{itemize}


% ──────────────────────────────────────────────────────────────
\subsection{Diferencias clave respecto a la aplicación móvil}
% ──────────────────────────────────────────────────────────────

\begin{table}[H]
  \centering
  \caption{Comparativa funcional breve: móvil vs web (módulos clave)}
  \label{tab:comparativa-web-movil}
  \begin{tabular}{p{3.2cm}p{4.8cm}p{4.8cm}}
    \toprule
    \textbf{Módulo} & \textbf{Móvil} & \textbf{Web} \\
    \midrule
    Escaneo ticket & Flujo centrado en cámara y múltiples capturas. & Flujo centrado en subida de archivo (PDF/imagen) + revisión en preview. \\
    Dashboard & Consulta rapida . & Énfasis en análisis visual y comprensión por período con gráficas más amplias. \\
    Reportes & ---. & Generación local PDF/Excel + gestión de historial con descarga y correo. \\
    Recomendaciones & ---. & Módulo dedicado con filtros mes/año y métricas calculadas por período. \\
    \bottomrule
  \end{tabular}
\end{table}


% ══════════════════════════════════════════════════════════════
\section{Backend}
% ══════════════════════════════════════════════════════════════

% Pendiente: integrar documentación del backend (Node.js + Express + MySQL).
% Incluir: stack técnico, arquitectura de la API REST, endpoints, base de
% datos, autenticación JWT, pipeline OCR y decisiones de diseño.

\begin{notabox}[Sección pendiente — Backend]
  Esta sección se encuentra pendiente de integración. La documentación
  técnica completa del backend de FINARA (Node.js + Express + MySQL) se
  encuentra en el documento paralelo \textit{``FINARA Backend ---
  Implementación Técnica''} y será incorporada en esta sección al
  consolidar el reporte final del Trabajo Terminal.
\end{notabox}
